% ============================================================================
% eskiz.tex -- el çizimi ("Excalidraw") görünümlü TikZ stilleri.
%
% Neden ayrı bir görsel dil: raporun geri kalanındaki şekiller
% (arka_plan.tex, deneyler.tex, gelecek_calismalar.tex) düz, keskin TikZ
% kutularıyla çizilmiş ve öyle kalıyor -- onlar ölçüm anlatıyor. Eğitim
% bölümü öğretmek için yazıldı; orada tahtaya çizilmiş gibi duran bir dil
% kullanılıyor, çünkü anlatılan şey bir sayı değil bir kurulum. İki dili
% karıştırmamak için eskiz stilleri burada, tek yerde duruyor.
%
% preamble.tex tarafından % ============================================================================
% eskiz.tex -- el çizimi ("Excalidraw") görünümlü TikZ stilleri.
%
% Neden ayrı bir görsel dil: raporun geri kalanındaki şekiller
% (arka_plan.tex, deneyler.tex, gelecek_calismalar.tex) düz, keskin TikZ
% kutularıyla çizilmiş ve öyle kalıyor -- onlar ölçüm anlatıyor. Eğitim
% bölümü öğretmek için yazıldı; orada tahtaya çizilmiş gibi duran bir dil
% kullanılıyor, çünkü anlatılan şey bir sayı değil bir kurulum. İki dili
% karıştırmamak için eskiz stilleri burada, tek yerde duruyor.
%
% preamble.tex tarafından % ============================================================================
% eskiz.tex -- el çizimi ("Excalidraw") görünümlü TikZ stilleri.
%
% Neden ayrı bir görsel dil: raporun geri kalanındaki şekiller
% (arka_plan.tex, deneyler.tex, gelecek_calismalar.tex) düz, keskin TikZ
% kutularıyla çizilmiş ve öyle kalıyor -- onlar ölçüm anlatıyor. Eğitim
% bölümü öğretmek için yazıldı; orada tahtaya çizilmiş gibi duran bir dil
% kullanılıyor, çünkü anlatılan şey bir sayı değil bir kurulum. İki dili
% karıştırmamak için eskiz stilleri burada, tek yerde duruyor.
%
% preamble.tex tarafından % ============================================================================
% eskiz.tex -- el çizimi ("Excalidraw") görünümlü TikZ stilleri.
%
% Neden ayrı bir görsel dil: raporun geri kalanındaki şekiller
% (arka_plan.tex, deneyler.tex, gelecek_calismalar.tex) düz, keskin TikZ
% kutularıyla çizilmiş ve öyle kalıyor -- onlar ölçüm anlatıyor. Eğitim
% bölümü öğretmek için yazıldı; orada tahtaya çizilmiş gibi duran bir dil
% kullanılıyor, çünkü anlatılan şey bir sayı değil bir kurulum. İki dili
% karıştırmamak için eskiz stilleri burada, tek yerde duruyor.
%
% preamble.tex tarafından \input{eskiz} ile çağrılır.
%
% Bir teknik not, yeniden kurcalayacak olan için: `rounded corners` ile
% `decorate` birlikte kullanılamıyor -- yuvarlatılmış köşenin minik yayları
% rastgele adım dekorasyonunda "Dimension too large" veriyor. Bu yüzden bütün
% kutular keskin köşeli; titremenin kendisi zaten köşeyi mekanik olmaktan
% çıkarıyor, ve dolgu ile kenar *aynı* dekore edilmiş yolu paylaştığı için
% dolgunun kenardan taşması diye bir sorun da doğmuyor.
% ============================================================================

\usetikzlibrary{decorations.pathmorphing, patterns, shapes.misc,
                shapes.geometric, calc, positioning, fit, backgrounds,
                arrows.meta}

% ---- Palet ------------------------------------------------------------------
% Excalidraw'ın kendi varsayılan paleti: her renk bir koyu "kalem" tonu ve ona
% eşlenen açık bir "dolgu" tonu. Raporun geri kalanının paletiyle
% (anarenk / vurgu / yesilvurgu) çakışmasın diye hepsi "es" önekli.
\definecolor{esmurekkep}{HTML}{1E1E1E}   % kalem: metin
\definecolor{escizgi}{HTML}{343A40}      % kalem: nötr kenar ve ok
\definecolor{esmavi}{HTML}{1971C2}       \definecolor{esmavid}{HTML}{A5D8FF}
\definecolor{esyesil}{HTML}{2F9E44}      \definecolor{esyesild}{HTML}{B2F2BB}
\definecolor{essari}{HTML}{E8590C}       \definecolor{essarid}{HTML}{FFEC99}
\definecolor{eskirmizi}{HTML}{E03131}    \definecolor{eskirmizid}{HTML}{FFC9C9}
\definecolor{esmor}{HTML}{6741D9}        \definecolor{esmord}{HTML}{D0BFFF}
\definecolor{esgri}{HTML}{868E96}        \definecolor{esgrid}{HTML}{E9ECEF}

% ---- Yazı tipi --------------------------------------------------------------
% Patrick Hand (SIL OFL 1.1, report/fonts/) -- Excalidraw'ın el yazısı fontuna
% (Virgil / Excalifont) en yakın duran, **tam Türkçe kapsamlı** serbest font.
% Depoda taşınıyor ki rapor hiçbir font indirmeden derlensin: Ubuntu'nun
% paketlediği el yazısı fontlarının ikisi de (Humor Sans, Comic Neue) ğ/Ğ/İ
% harflerini içermiyor, yani Türkçe bir şekilde kullanılamıyorlar.
%
% Dosya yoksa \elyazi sessizce TeX Gyre Heros'a düşer: şekiller aynı çizilir,
% yalnızca etiketler düz sans olur ve derleme kırılmaz.
%
% AutoFakeBold / AutoFakeSlant: Patrick Hand tek ağırlıkta geliyor, kalın ve
% italik varyantı yok. Şekillerde \textbf ve \emph kullanılıyor (bir kutunun
% adı kalın, bir yan not italik), o yüzden ikisi de sentetik üretiliyor --
% aksi hâlde XeLaTeX her kullanımda "font shape undefined" uyarısı verip
% sessizce düz varyanta düşerdi ve vurgu kaybolurdu.
%
% Ligatures=TeX: `---` ve `--` girdilerini uzun/kısa tireye çeviren klasik
% TeX bağı. \newfontfamily ile kurulan bir aileye kendiliğinden gelmiyor,
% o yüzden açıkça isteniyor -- yoksa şekillerdeki her uzun tire üç ayrı
% kısa çizgi olarak dizilir (fontun kendisinde iki glif de var).
\IfFileExists{fonts/PatrickHand-Regular.ttf}{%
  \newfontfamily\elyazi{PatrickHand-Regular.ttf}[%
    Path=fonts/, Scale=1.18, AutoFakeBold=1.9, AutoFakeSlant=0.2,
    Ligatures=TeX]%
}{%
  \newcommand{\elyazi}{\sffamily}%
}

% ---- Çizgi: elle çizilmiş titreme -------------------------------------------
% `random steps` yolu küçük rastgele adımlara bölüyor. Genlik ve adım boyu
% deneyerek seçildi: 0,7 pt / 2,6 mm elle çizilmiş okunuyor; daha büyüğü
% suda dalgalanmış gibi duruyor, daha küçüğü düz cetvel.
%
% Tohum her şeklin başında sabitleniyor (\eskizbasla), böylece aynı kaynak her
% derlemede baytı baytına aynı PDF'i veriyor -- şekiller derlemeden derlemeye
% "kıpırdamıyor".
% Titreme ayarları iki adlandırılmış stilde duruyor. \newcommand ile bir
% makroya koyup `decoration={random steps, \makro}` yazmak *çalışmıyor* --
% pgfkeys anahtar listesini genişletmeden okuyor ve makronun adını anahtar
% sanıyor ("I do not know the key '/pgf/decoration/.expanded'"). Stil olarak
% tanımlamak hem çalışıyor hem de tek yerden ayarlanabilir kalıyor.
\tikzset{
  estitreme/.style={decorate,
    decoration={random steps, segment length=2.6mm, amplitude=0.7pt}},
  estitremeok/.style={decorate,
    decoration={random steps, segment length=4mm, amplitude=0.9pt}},
}

\tikzset{
  cizik/.style={
    estitreme,
    line width=1pt, line cap=round, line join=round,
  },
}

% Excalidraw her kenarı iki kez, hafifçe kaydırarak çiziyor. İkinci geçiş
% burada bir `postaction`: aynı yol yeniden dekore edilince rastgele adımlar
% yeniden çekiliyor, yani ikinci çizgi birincisinin tam üstüne oturmuyor --
% aranan etki tam olarak bu.
\tikzset{
  eskutu/.style n args={2}{
    cizik, draw=#1, fill=#2,
    align=center, inner sep=2.2mm, text=esmurekkep, font=\elyazi\small,
    postaction={
      estitreme,
      draw=#1, line width=0.55pt, opacity=0.45,
      line cap=round, line join=round,
    },
  },
}

% Renkli kutular. `kbos` dolgusuz (kağıdın kendisi), gerisi Excalidraw'ın
% açık dolgularıyla. Bu bölümdeki anlam sözleşmesi:
%   mavi   -> eğitilen modül        yeşil -> hedef / doğru sonuç
%   sarı   -> veri, ara tensör      kırmızı -> hata, risk, kaybedilen şey
%   mor    -> öğretmen / dış model  gri   -> nötr, bilgi
\tikzset{
  kbos/.style={eskutu={escizgi}{white}},
  kmavi/.style={eskutu={esmavi}{esmavid}},
  kyesil/.style={eskutu={esyesil}{esyesild}},
  ksari/.style={eskutu={essari}{essarid}},
  kkirmizi/.style={eskutu={eskirmizi}{eskirmizid}},
  kmor/.style={eskutu={esmor}{esmord}},
  kgri/.style={eskutu={esgri}{esgrid}},
}

% Donmuş (eğitilmeyen) modül: taralı dolgu -- Excalidraw'ın "hachure"
% doldurması. Bu bölümde "bu kutuya dokunulmuyor" demenin görsel yolu, ve
% bilerek renksiz: donuk olmak bir renk değil, bir yokluk.
\tikzset{
  kdonuk/.style={
    cizik, draw=esgri,
    pattern=north east lines, pattern color=esgri!50,
    align=center, inner sep=2.2mm, text=esgri, font=\elyazi\small,
  },
}

% ---- Oklar ------------------------------------------------------------------
% `Straight Barb` = açık V uç; Excalidraw'ın ok ucu bu.
\tikzset{
  esokgovde/.style={
    estitremeok,
    line cap=round, line join=round,
  },
  esok/.style={esokgovde, draw=escizgi, line width=1pt,
               -{Straight Barb[length=2mm, width=3mm, line width=1pt]}},
  % renkli ok: \draw[esokr=esmavi] ...
  esokr/.style={esokgovde, draw=#1, line width=1pt,
                -{Straight Barb[length=2mm, width=3mm, line width=1pt]}},
  % ince kesikli: "veri akışı değil, ilişki" oku
  esink/.style={esokgovde, draw=esgri, line width=0.7pt,
                dash pattern=on 1.7mm off 1.2mm,
                -{Straight Barb[length=1.6mm, width=2.4mm, line width=0.7pt]}},
  % uçsuz çizgi
  escizik/.style={esokgovde, draw=escizgi, line width=1pt},
  escizikr/.style={esokgovde, draw=#1, line width=1pt},
}

% ---- Etiketler ---------------------------------------------------------------
\tikzset{
  esnot/.style={font=\elyazi\footnotesize, text=esgri, align=center, inner sep=1.5pt},
  esvurgu/.style={font=\elyazi\footnotesize, text=eskirmizi, align=center, inner sep=1.5pt},
  esiyi/.style={font=\elyazi\footnotesize, text=esyesil, align=center, inner sep=1.5pt},
  esbaslik/.style={font=\elyazi\small, text=esmurekkep, align=center, inner sep=1.5pt},
}

% Elle çizilmiş bir çerçeve içine alma -- bir grubu "işaretlemek" için.
% Kullanım: \node[esisaret, fit=(a)(b)] {};   ya da başlıklısı için altına
% ayrı bir \node[esnot] koyun.
\tikzset{
  esisaret/.style={
    cizik, draw=eskirmizi, line width=1pt, inner sep=2.8mm,
    postaction={estitreme,
                draw=eskirmizi, line width=0.55pt, opacity=0.45,
                line cap=round, line join=round},
  },
  esisaretr/.style={
    cizik, draw=#1, line width=1pt, inner sep=2.8mm,
    postaction={estitreme,
                draw=#1, line width=0.55pt, opacity=0.45,
                line cap=round, line join=round},
  },
}

% ---- Şekil başlangıcı ---------------------------------------------------------
% Her el çizimi şeklin başına konur: rastgele tohumu sabitler (derlemeler
% arası kararlılık). İsteğe bağlı argüman tohumu değiştirir, yani bir şeklin
% titremesi beğenilmezse sadece sayıyı değiştirmek yetiyor.
\newcommand{\eskizbasla}[1][2024]{\pgfmathsetseed{#1}}
 ile çağrılır.
%
% Bir teknik not, yeniden kurcalayacak olan için: `rounded corners` ile
% `decorate` birlikte kullanılamıyor -- yuvarlatılmış köşenin minik yayları
% rastgele adım dekorasyonunda "Dimension too large" veriyor. Bu yüzden bütün
% kutular keskin köşeli; titremenin kendisi zaten köşeyi mekanik olmaktan
% çıkarıyor, ve dolgu ile kenar *aynı* dekore edilmiş yolu paylaştığı için
% dolgunun kenardan taşması diye bir sorun da doğmuyor.
% ============================================================================

\usetikzlibrary{decorations.pathmorphing, patterns, shapes.misc,
                shapes.geometric, calc, positioning, fit, backgrounds,
                arrows.meta}

% ---- Palet ------------------------------------------------------------------
% Excalidraw'ın kendi varsayılan paleti: her renk bir koyu "kalem" tonu ve ona
% eşlenen açık bir "dolgu" tonu. Raporun geri kalanının paletiyle
% (anarenk / vurgu / yesilvurgu) çakışmasın diye hepsi "es" önekli.
\definecolor{esmurekkep}{HTML}{1E1E1E}   % kalem: metin
\definecolor{escizgi}{HTML}{343A40}      % kalem: nötr kenar ve ok
\definecolor{esmavi}{HTML}{1971C2}       \definecolor{esmavid}{HTML}{A5D8FF}
\definecolor{esyesil}{HTML}{2F9E44}      \definecolor{esyesild}{HTML}{B2F2BB}
\definecolor{essari}{HTML}{E8590C}       \definecolor{essarid}{HTML}{FFEC99}
\definecolor{eskirmizi}{HTML}{E03131}    \definecolor{eskirmizid}{HTML}{FFC9C9}
\definecolor{esmor}{HTML}{6741D9}        \definecolor{esmord}{HTML}{D0BFFF}
\definecolor{esgri}{HTML}{868E96}        \definecolor{esgrid}{HTML}{E9ECEF}

% ---- Yazı tipi --------------------------------------------------------------
% Patrick Hand (SIL OFL 1.1, report/fonts/) -- Excalidraw'ın el yazısı fontuna
% (Virgil / Excalifont) en yakın duran, **tam Türkçe kapsamlı** serbest font.
% Depoda taşınıyor ki rapor hiçbir font indirmeden derlensin: Ubuntu'nun
% paketlediği el yazısı fontlarının ikisi de (Humor Sans, Comic Neue) ğ/Ğ/İ
% harflerini içermiyor, yani Türkçe bir şekilde kullanılamıyorlar.
%
% Dosya yoksa \elyazi sessizce TeX Gyre Heros'a düşer: şekiller aynı çizilir,
% yalnızca etiketler düz sans olur ve derleme kırılmaz.
%
% AutoFakeBold / AutoFakeSlant: Patrick Hand tek ağırlıkta geliyor, kalın ve
% italik varyantı yok. Şekillerde \textbf ve \emph kullanılıyor (bir kutunun
% adı kalın, bir yan not italik), o yüzden ikisi de sentetik üretiliyor --
% aksi hâlde XeLaTeX her kullanımda "font shape undefined" uyarısı verip
% sessizce düz varyanta düşerdi ve vurgu kaybolurdu.
%
% Ligatures=TeX: `---` ve `--` girdilerini uzun/kısa tireye çeviren klasik
% TeX bağı. \newfontfamily ile kurulan bir aileye kendiliğinden gelmiyor,
% o yüzden açıkça isteniyor -- yoksa şekillerdeki her uzun tire üç ayrı
% kısa çizgi olarak dizilir (fontun kendisinde iki glif de var).
\IfFileExists{fonts/PatrickHand-Regular.ttf}{%
  \newfontfamily\elyazi{PatrickHand-Regular.ttf}[%
    Path=fonts/, Scale=1.18, AutoFakeBold=1.9, AutoFakeSlant=0.2,
    Ligatures=TeX]%
}{%
  \newcommand{\elyazi}{\sffamily}%
}

% ---- Çizgi: elle çizilmiş titreme -------------------------------------------
% `random steps` yolu küçük rastgele adımlara bölüyor. Genlik ve adım boyu
% deneyerek seçildi: 0,7 pt / 2,6 mm elle çizilmiş okunuyor; daha büyüğü
% suda dalgalanmış gibi duruyor, daha küçüğü düz cetvel.
%
% Tohum her şeklin başında sabitleniyor (\eskizbasla), böylece aynı kaynak her
% derlemede baytı baytına aynı PDF'i veriyor -- şekiller derlemeden derlemeye
% "kıpırdamıyor".
% Titreme ayarları iki adlandırılmış stilde duruyor. \newcommand ile bir
% makroya koyup `decoration={random steps, \makro}` yazmak *çalışmıyor* --
% pgfkeys anahtar listesini genişletmeden okuyor ve makronun adını anahtar
% sanıyor ("I do not know the key '/pgf/decoration/.expanded'"). Stil olarak
% tanımlamak hem çalışıyor hem de tek yerden ayarlanabilir kalıyor.
\tikzset{
  estitreme/.style={decorate,
    decoration={random steps, segment length=2.6mm, amplitude=0.7pt}},
  estitremeok/.style={decorate,
    decoration={random steps, segment length=4mm, amplitude=0.9pt}},
}

\tikzset{
  cizik/.style={
    estitreme,
    line width=1pt, line cap=round, line join=round,
  },
}

% Excalidraw her kenarı iki kez, hafifçe kaydırarak çiziyor. İkinci geçiş
% burada bir `postaction`: aynı yol yeniden dekore edilince rastgele adımlar
% yeniden çekiliyor, yani ikinci çizgi birincisinin tam üstüne oturmuyor --
% aranan etki tam olarak bu.
\tikzset{
  eskutu/.style n args={2}{
    cizik, draw=#1, fill=#2,
    align=center, inner sep=2.2mm, text=esmurekkep, font=\elyazi\small,
    postaction={
      estitreme,
      draw=#1, line width=0.55pt, opacity=0.45,
      line cap=round, line join=round,
    },
  },
}

% Renkli kutular. `kbos` dolgusuz (kağıdın kendisi), gerisi Excalidraw'ın
% açık dolgularıyla. Bu bölümdeki anlam sözleşmesi:
%   mavi   -> eğitilen modül        yeşil -> hedef / doğru sonuç
%   sarı   -> veri, ara tensör      kırmızı -> hata, risk, kaybedilen şey
%   mor    -> öğretmen / dış model  gri   -> nötr, bilgi
\tikzset{
  kbos/.style={eskutu={escizgi}{white}},
  kmavi/.style={eskutu={esmavi}{esmavid}},
  kyesil/.style={eskutu={esyesil}{esyesild}},
  ksari/.style={eskutu={essari}{essarid}},
  kkirmizi/.style={eskutu={eskirmizi}{eskirmizid}},
  kmor/.style={eskutu={esmor}{esmord}},
  kgri/.style={eskutu={esgri}{esgrid}},
}

% Donmuş (eğitilmeyen) modül: taralı dolgu -- Excalidraw'ın "hachure"
% doldurması. Bu bölümde "bu kutuya dokunulmuyor" demenin görsel yolu, ve
% bilerek renksiz: donuk olmak bir renk değil, bir yokluk.
\tikzset{
  kdonuk/.style={
    cizik, draw=esgri,
    pattern=north east lines, pattern color=esgri!50,
    align=center, inner sep=2.2mm, text=esgri, font=\elyazi\small,
  },
}

% ---- Oklar ------------------------------------------------------------------
% `Straight Barb` = açık V uç; Excalidraw'ın ok ucu bu.
\tikzset{
  esokgovde/.style={
    estitremeok,
    line cap=round, line join=round,
  },
  esok/.style={esokgovde, draw=escizgi, line width=1pt,
               -{Straight Barb[length=2mm, width=3mm, line width=1pt]}},
  % renkli ok: \draw[esokr=esmavi] ...
  esokr/.style={esokgovde, draw=#1, line width=1pt,
                -{Straight Barb[length=2mm, width=3mm, line width=1pt]}},
  % ince kesikli: "veri akışı değil, ilişki" oku
  esink/.style={esokgovde, draw=esgri, line width=0.7pt,
                dash pattern=on 1.7mm off 1.2mm,
                -{Straight Barb[length=1.6mm, width=2.4mm, line width=0.7pt]}},
  % uçsuz çizgi
  escizik/.style={esokgovde, draw=escizgi, line width=1pt},
  escizikr/.style={esokgovde, draw=#1, line width=1pt},
}

% ---- Etiketler ---------------------------------------------------------------
\tikzset{
  esnot/.style={font=\elyazi\footnotesize, text=esgri, align=center, inner sep=1.5pt},
  esvurgu/.style={font=\elyazi\footnotesize, text=eskirmizi, align=center, inner sep=1.5pt},
  esiyi/.style={font=\elyazi\footnotesize, text=esyesil, align=center, inner sep=1.5pt},
  esbaslik/.style={font=\elyazi\small, text=esmurekkep, align=center, inner sep=1.5pt},
}

% Elle çizilmiş bir çerçeve içine alma -- bir grubu "işaretlemek" için.
% Kullanım: \node[esisaret, fit=(a)(b)] {};   ya da başlıklısı için altına
% ayrı bir \node[esnot] koyun.
\tikzset{
  esisaret/.style={
    cizik, draw=eskirmizi, line width=1pt, inner sep=2.8mm,
    postaction={estitreme,
                draw=eskirmizi, line width=0.55pt, opacity=0.45,
                line cap=round, line join=round},
  },
  esisaretr/.style={
    cizik, draw=#1, line width=1pt, inner sep=2.8mm,
    postaction={estitreme,
                draw=#1, line width=0.55pt, opacity=0.45,
                line cap=round, line join=round},
  },
}

% ---- Şekil başlangıcı ---------------------------------------------------------
% Her el çizimi şeklin başına konur: rastgele tohumu sabitler (derlemeler
% arası kararlılık). İsteğe bağlı argüman tohumu değiştirir, yani bir şeklin
% titremesi beğenilmezse sadece sayıyı değiştirmek yetiyor.
\newcommand{\eskizbasla}[1][2024]{\pgfmathsetseed{#1}}
 ile çağrılır.
%
% Bir teknik not, yeniden kurcalayacak olan için: `rounded corners` ile
% `decorate` birlikte kullanılamıyor -- yuvarlatılmış köşenin minik yayları
% rastgele adım dekorasyonunda "Dimension too large" veriyor. Bu yüzden bütün
% kutular keskin köşeli; titremenin kendisi zaten köşeyi mekanik olmaktan
% çıkarıyor, ve dolgu ile kenar *aynı* dekore edilmiş yolu paylaştığı için
% dolgunun kenardan taşması diye bir sorun da doğmuyor.
% ============================================================================

\usetikzlibrary{decorations.pathmorphing, patterns, shapes.misc,
                shapes.geometric, calc, positioning, fit, backgrounds,
                arrows.meta}

% ---- Palet ------------------------------------------------------------------
% Excalidraw'ın kendi varsayılan paleti: her renk bir koyu "kalem" tonu ve ona
% eşlenen açık bir "dolgu" tonu. Raporun geri kalanının paletiyle
% (anarenk / vurgu / yesilvurgu) çakışmasın diye hepsi "es" önekli.
\definecolor{esmurekkep}{HTML}{1E1E1E}   % kalem: metin
\definecolor{escizgi}{HTML}{343A40}      % kalem: nötr kenar ve ok
\definecolor{esmavi}{HTML}{1971C2}       \definecolor{esmavid}{HTML}{A5D8FF}
\definecolor{esyesil}{HTML}{2F9E44}      \definecolor{esyesild}{HTML}{B2F2BB}
\definecolor{essari}{HTML}{E8590C}       \definecolor{essarid}{HTML}{FFEC99}
\definecolor{eskirmizi}{HTML}{E03131}    \definecolor{eskirmizid}{HTML}{FFC9C9}
\definecolor{esmor}{HTML}{6741D9}        \definecolor{esmord}{HTML}{D0BFFF}
\definecolor{esgri}{HTML}{868E96}        \definecolor{esgrid}{HTML}{E9ECEF}

% ---- Yazı tipi --------------------------------------------------------------
% Patrick Hand (SIL OFL 1.1, report/fonts/) -- Excalidraw'ın el yazısı fontuna
% (Virgil / Excalifont) en yakın duran, **tam Türkçe kapsamlı** serbest font.
% Depoda taşınıyor ki rapor hiçbir font indirmeden derlensin: Ubuntu'nun
% paketlediği el yazısı fontlarının ikisi de (Humor Sans, Comic Neue) ğ/Ğ/İ
% harflerini içermiyor, yani Türkçe bir şekilde kullanılamıyorlar.
%
% Dosya yoksa \elyazi sessizce TeX Gyre Heros'a düşer: şekiller aynı çizilir,
% yalnızca etiketler düz sans olur ve derleme kırılmaz.
%
% AutoFakeBold / AutoFakeSlant: Patrick Hand tek ağırlıkta geliyor, kalın ve
% italik varyantı yok. Şekillerde \textbf ve \emph kullanılıyor (bir kutunun
% adı kalın, bir yan not italik), o yüzden ikisi de sentetik üretiliyor --
% aksi hâlde XeLaTeX her kullanımda "font shape undefined" uyarısı verip
% sessizce düz varyanta düşerdi ve vurgu kaybolurdu.
%
% Ligatures=TeX: `---` ve `--` girdilerini uzun/kısa tireye çeviren klasik
% TeX bağı. \newfontfamily ile kurulan bir aileye kendiliğinden gelmiyor,
% o yüzden açıkça isteniyor -- yoksa şekillerdeki her uzun tire üç ayrı
% kısa çizgi olarak dizilir (fontun kendisinde iki glif de var).
\IfFileExists{fonts/PatrickHand-Regular.ttf}{%
  \newfontfamily\elyazi{PatrickHand-Regular.ttf}[%
    Path=fonts/, Scale=1.18, AutoFakeBold=1.9, AutoFakeSlant=0.2,
    Ligatures=TeX]%
}{%
  \newcommand{\elyazi}{\sffamily}%
}

% ---- Çizgi: elle çizilmiş titreme -------------------------------------------
% `random steps` yolu küçük rastgele adımlara bölüyor. Genlik ve adım boyu
% deneyerek seçildi: 0,7 pt / 2,6 mm elle çizilmiş okunuyor; daha büyüğü
% suda dalgalanmış gibi duruyor, daha küçüğü düz cetvel.
%
% Tohum her şeklin başında sabitleniyor (\eskizbasla), böylece aynı kaynak her
% derlemede baytı baytına aynı PDF'i veriyor -- şekiller derlemeden derlemeye
% "kıpırdamıyor".
% Titreme ayarları iki adlandırılmış stilde duruyor. \newcommand ile bir
% makroya koyup `decoration={random steps, \makro}` yazmak *çalışmıyor* --
% pgfkeys anahtar listesini genişletmeden okuyor ve makronun adını anahtar
% sanıyor ("I do not know the key '/pgf/decoration/.expanded'"). Stil olarak
% tanımlamak hem çalışıyor hem de tek yerden ayarlanabilir kalıyor.
\tikzset{
  estitreme/.style={decorate,
    decoration={random steps, segment length=2.6mm, amplitude=0.7pt}},
  estitremeok/.style={decorate,
    decoration={random steps, segment length=4mm, amplitude=0.9pt}},
}

\tikzset{
  cizik/.style={
    estitreme,
    line width=1pt, line cap=round, line join=round,
  },
}

% Excalidraw her kenarı iki kez, hafifçe kaydırarak çiziyor. İkinci geçiş
% burada bir `postaction`: aynı yol yeniden dekore edilince rastgele adımlar
% yeniden çekiliyor, yani ikinci çizgi birincisinin tam üstüne oturmuyor --
% aranan etki tam olarak bu.
\tikzset{
  eskutu/.style n args={2}{
    cizik, draw=#1, fill=#2,
    align=center, inner sep=2.2mm, text=esmurekkep, font=\elyazi\small,
    postaction={
      estitreme,
      draw=#1, line width=0.55pt, opacity=0.45,
      line cap=round, line join=round,
    },
  },
}

% Renkli kutular. `kbos` dolgusuz (kağıdın kendisi), gerisi Excalidraw'ın
% açık dolgularıyla. Bu bölümdeki anlam sözleşmesi:
%   mavi   -> eğitilen modül        yeşil -> hedef / doğru sonuç
%   sarı   -> veri, ara tensör      kırmızı -> hata, risk, kaybedilen şey
%   mor    -> öğretmen / dış model  gri   -> nötr, bilgi
\tikzset{
  kbos/.style={eskutu={escizgi}{white}},
  kmavi/.style={eskutu={esmavi}{esmavid}},
  kyesil/.style={eskutu={esyesil}{esyesild}},
  ksari/.style={eskutu={essari}{essarid}},
  kkirmizi/.style={eskutu={eskirmizi}{eskirmizid}},
  kmor/.style={eskutu={esmor}{esmord}},
  kgri/.style={eskutu={esgri}{esgrid}},
}

% Donmuş (eğitilmeyen) modül: taralı dolgu -- Excalidraw'ın "hachure"
% doldurması. Bu bölümde "bu kutuya dokunulmuyor" demenin görsel yolu, ve
% bilerek renksiz: donuk olmak bir renk değil, bir yokluk.
\tikzset{
  kdonuk/.style={
    cizik, draw=esgri,
    pattern=north east lines, pattern color=esgri!50,
    align=center, inner sep=2.2mm, text=esgri, font=\elyazi\small,
  },
}

% ---- Oklar ------------------------------------------------------------------
% `Straight Barb` = açık V uç; Excalidraw'ın ok ucu bu.
\tikzset{
  esokgovde/.style={
    estitremeok,
    line cap=round, line join=round,
  },
  esok/.style={esokgovde, draw=escizgi, line width=1pt,
               -{Straight Barb[length=2mm, width=3mm, line width=1pt]}},
  % renkli ok: \draw[esokr=esmavi] ...
  esokr/.style={esokgovde, draw=#1, line width=1pt,
                -{Straight Barb[length=2mm, width=3mm, line width=1pt]}},
  % ince kesikli: "veri akışı değil, ilişki" oku
  esink/.style={esokgovde, draw=esgri, line width=0.7pt,
                dash pattern=on 1.7mm off 1.2mm,
                -{Straight Barb[length=1.6mm, width=2.4mm, line width=0.7pt]}},
  % uçsuz çizgi
  escizik/.style={esokgovde, draw=escizgi, line width=1pt},
  escizikr/.style={esokgovde, draw=#1, line width=1pt},
}

% ---- Etiketler ---------------------------------------------------------------
\tikzset{
  esnot/.style={font=\elyazi\footnotesize, text=esgri, align=center, inner sep=1.5pt},
  esvurgu/.style={font=\elyazi\footnotesize, text=eskirmizi, align=center, inner sep=1.5pt},
  esiyi/.style={font=\elyazi\footnotesize, text=esyesil, align=center, inner sep=1.5pt},
  esbaslik/.style={font=\elyazi\small, text=esmurekkep, align=center, inner sep=1.5pt},
}

% Elle çizilmiş bir çerçeve içine alma -- bir grubu "işaretlemek" için.
% Kullanım: \node[esisaret, fit=(a)(b)] {};   ya da başlıklısı için altına
% ayrı bir \node[esnot] koyun.
\tikzset{
  esisaret/.style={
    cizik, draw=eskirmizi, line width=1pt, inner sep=2.8mm,
    postaction={estitreme,
                draw=eskirmizi, line width=0.55pt, opacity=0.45,
                line cap=round, line join=round},
  },
  esisaretr/.style={
    cizik, draw=#1, line width=1pt, inner sep=2.8mm,
    postaction={estitreme,
                draw=#1, line width=0.55pt, opacity=0.45,
                line cap=round, line join=round},
  },
}

% ---- Şekil başlangıcı ---------------------------------------------------------
% Her el çizimi şeklin başına konur: rastgele tohumu sabitler (derlemeler
% arası kararlılık). İsteğe bağlı argüman tohumu değiştirir, yani bir şeklin
% titremesi beğenilmezse sadece sayıyı değiştirmek yetiyor.
\newcommand{\eskizbasla}[1][2024]{\pgfmathsetseed{#1}}
 ile çağrılır.
%
% Bir teknik not, yeniden kurcalayacak olan için: `rounded corners` ile
% `decorate` birlikte kullanılamıyor -- yuvarlatılmış köşenin minik yayları
% rastgele adım dekorasyonunda "Dimension too large" veriyor. Bu yüzden bütün
% kutular keskin köşeli; titremenin kendisi zaten köşeyi mekanik olmaktan
% çıkarıyor, ve dolgu ile kenar *aynı* dekore edilmiş yolu paylaştığı için
% dolgunun kenardan taşması diye bir sorun da doğmuyor.
% ============================================================================

\usetikzlibrary{decorations.pathmorphing, patterns, shapes.misc,
                shapes.geometric, calc, positioning, fit, backgrounds,
                arrows.meta}

% ---- Palet ------------------------------------------------------------------
% Excalidraw'ın kendi varsayılan paleti: her renk bir koyu "kalem" tonu ve ona
% eşlenen açık bir "dolgu" tonu. Raporun geri kalanının paletiyle
% (anarenk / vurgu / yesilvurgu) çakışmasın diye hepsi "es" önekli.
\definecolor{esmurekkep}{HTML}{1E1E1E}   % kalem: metin
\definecolor{escizgi}{HTML}{343A40}      % kalem: nötr kenar ve ok
\definecolor{esmavi}{HTML}{1971C2}       \definecolor{esmavid}{HTML}{A5D8FF}
\definecolor{esyesil}{HTML}{2F9E44}      \definecolor{esyesild}{HTML}{B2F2BB}
\definecolor{essari}{HTML}{E8590C}       \definecolor{essarid}{HTML}{FFEC99}
\definecolor{eskirmizi}{HTML}{E03131}    \definecolor{eskirmizid}{HTML}{FFC9C9}
\definecolor{esmor}{HTML}{6741D9}        \definecolor{esmord}{HTML}{D0BFFF}
\definecolor{esgri}{HTML}{868E96}        \definecolor{esgrid}{HTML}{E9ECEF}

% ---- Yazı tipi --------------------------------------------------------------
% Patrick Hand (SIL OFL 1.1, report/fonts/) -- Excalidraw'ın el yazısı fontuna
% (Virgil / Excalifont) en yakın duran, **tam Türkçe kapsamlı** serbest font.
% Depoda taşınıyor ki rapor hiçbir font indirmeden derlensin: Ubuntu'nun
% paketlediği el yazısı fontlarının ikisi de (Humor Sans, Comic Neue) ğ/Ğ/İ
% harflerini içermiyor, yani Türkçe bir şekilde kullanılamıyorlar.
%
% Dosya yoksa \elyazi sessizce TeX Gyre Heros'a düşer: şekiller aynı çizilir,
% yalnızca etiketler düz sans olur ve derleme kırılmaz.
%
% AutoFakeBold / AutoFakeSlant: Patrick Hand tek ağırlıkta geliyor, kalın ve
% italik varyantı yok. Şekillerde \textbf ve \emph kullanılıyor (bir kutunun
% adı kalın, bir yan not italik), o yüzden ikisi de sentetik üretiliyor --
% aksi hâlde XeLaTeX her kullanımda "font shape undefined" uyarısı verip
% sessizce düz varyanta düşerdi ve vurgu kaybolurdu.
%
% Ligatures=TeX: `---` ve `--` girdilerini uzun/kısa tireye çeviren klasik
% TeX bağı. \newfontfamily ile kurulan bir aileye kendiliğinden gelmiyor,
% o yüzden açıkça isteniyor -- yoksa şekillerdeki her uzun tire üç ayrı
% kısa çizgi olarak dizilir (fontun kendisinde iki glif de var).
\IfFileExists{fonts/PatrickHand-Regular.ttf}{%
  \newfontfamily\elyazi{PatrickHand-Regular.ttf}[%
    Path=fonts/, Scale=1.18, AutoFakeBold=1.9, AutoFakeSlant=0.2,
    Ligatures=TeX]%
}{%
  \newcommand{\elyazi}{\sffamily}%
}

% ---- Çizgi: elle çizilmiş titreme -------------------------------------------
% `random steps` yolu küçük rastgele adımlara bölüyor. Genlik ve adım boyu
% deneyerek seçildi: 0,7 pt / 2,6 mm elle çizilmiş okunuyor; daha büyüğü
% suda dalgalanmış gibi duruyor, daha küçüğü düz cetvel.
%
% Tohum her şeklin başında sabitleniyor (\eskizbasla), böylece aynı kaynak her
% derlemede baytı baytına aynı PDF'i veriyor -- şekiller derlemeden derlemeye
% "kıpırdamıyor".
% Titreme ayarları iki adlandırılmış stilde duruyor. \newcommand ile bir
% makroya koyup `decoration={random steps, \makro}` yazmak *çalışmıyor* --
% pgfkeys anahtar listesini genişletmeden okuyor ve makronun adını anahtar
% sanıyor ("I do not know the key '/pgf/decoration/.expanded'"). Stil olarak
% tanımlamak hem çalışıyor hem de tek yerden ayarlanabilir kalıyor.
\tikzset{
  estitreme/.style={decorate,
    decoration={random steps, segment length=2.6mm, amplitude=0.7pt}},
  estitremeok/.style={decorate,
    decoration={random steps, segment length=4mm, amplitude=0.9pt}},
}

\tikzset{
  cizik/.style={
    estitreme,
    line width=1pt, line cap=round, line join=round,
  },
}

% Excalidraw her kenarı iki kez, hafifçe kaydırarak çiziyor. İkinci geçiş
% burada bir `postaction`: aynı yol yeniden dekore edilince rastgele adımlar
% yeniden çekiliyor, yani ikinci çizgi birincisinin tam üstüne oturmuyor --
% aranan etki tam olarak bu.
\tikzset{
  eskutu/.style n args={2}{
    cizik, draw=#1, fill=#2,
    align=center, inner sep=2.2mm, text=esmurekkep, font=\elyazi\small,
    postaction={
      estitreme,
      draw=#1, line width=0.55pt, opacity=0.45,
      line cap=round, line join=round,
    },
  },
}

% Renkli kutular. `kbos` dolgusuz (kağıdın kendisi), gerisi Excalidraw'ın
% açık dolgularıyla. Bu bölümdeki anlam sözleşmesi:
%   mavi   -> eğitilen modül        yeşil -> hedef / doğru sonuç
%   sarı   -> veri, ara tensör      kırmızı -> hata, risk, kaybedilen şey
%   mor    -> öğretmen / dış model  gri   -> nötr, bilgi
\tikzset{
  kbos/.style={eskutu={escizgi}{white}},
  kmavi/.style={eskutu={esmavi}{esmavid}},
  kyesil/.style={eskutu={esyesil}{esyesild}},
  ksari/.style={eskutu={essari}{essarid}},
  kkirmizi/.style={eskutu={eskirmizi}{eskirmizid}},
  kmor/.style={eskutu={esmor}{esmord}},
  kgri/.style={eskutu={esgri}{esgrid}},
}

% Donmuş (eğitilmeyen) modül: taralı dolgu -- Excalidraw'ın "hachure"
% doldurması. Bu bölümde "bu kutuya dokunulmuyor" demenin görsel yolu, ve
% bilerek renksiz: donuk olmak bir renk değil, bir yokluk.
\tikzset{
  kdonuk/.style={
    cizik, draw=esgri,
    pattern=north east lines, pattern color=esgri!50,
    align=center, inner sep=2.2mm, text=esgri, font=\elyazi\small,
  },
}

% ---- Oklar ------------------------------------------------------------------
% `Straight Barb` = açık V uç; Excalidraw'ın ok ucu bu.
\tikzset{
  esokgovde/.style={
    estitremeok,
    line cap=round, line join=round,
  },
  esok/.style={esokgovde, draw=escizgi, line width=1pt,
               -{Straight Barb[length=2mm, width=3mm, line width=1pt]}},
  % renkli ok: \draw[esokr=esmavi] ...
  esokr/.style={esokgovde, draw=#1, line width=1pt,
                -{Straight Barb[length=2mm, width=3mm, line width=1pt]}},
  % ince kesikli: "veri akışı değil, ilişki" oku
  esink/.style={esokgovde, draw=esgri, line width=0.7pt,
                dash pattern=on 1.7mm off 1.2mm,
                -{Straight Barb[length=1.6mm, width=2.4mm, line width=0.7pt]}},
  % uçsuz çizgi
  escizik/.style={esokgovde, draw=escizgi, line width=1pt},
  escizikr/.style={esokgovde, draw=#1, line width=1pt},
}

% ---- Etiketler ---------------------------------------------------------------
\tikzset{
  esnot/.style={font=\elyazi\footnotesize, text=esgri, align=center, inner sep=1.5pt},
  esvurgu/.style={font=\elyazi\footnotesize, text=eskirmizi, align=center, inner sep=1.5pt},
  esiyi/.style={font=\elyazi\footnotesize, text=esyesil, align=center, inner sep=1.5pt},
  esbaslik/.style={font=\elyazi\small, text=esmurekkep, align=center, inner sep=1.5pt},
}

% Elle çizilmiş bir çerçeve içine alma -- bir grubu "işaretlemek" için.
% Kullanım: \node[esisaret, fit=(a)(b)] {};   ya da başlıklısı için altına
% ayrı bir \node[esnot] koyun.
\tikzset{
  esisaret/.style={
    cizik, draw=eskirmizi, line width=1pt, inner sep=2.8mm,
    postaction={estitreme,
                draw=eskirmizi, line width=0.55pt, opacity=0.45,
                line cap=round, line join=round},
  },
  esisaretr/.style={
    cizik, draw=#1, line width=1pt, inner sep=2.8mm,
    postaction={estitreme,
                draw=#1, line width=0.55pt, opacity=0.45,
                line cap=round, line join=round},
  },
}

% ---- Şekil başlangıcı ---------------------------------------------------------
% Her el çizimi şeklin başına konur: rastgele tohumu sabitler (derlemeler
% arası kararlılık). İsteğe bağlı argüman tohumu değiştirir, yani bir şeklin
% titremesi beğenilmezse sadece sayıyı değiştirmek yetiyor.
\newcommand{\eskizbasla}[1][2024]{\pgfmathsetseed{#1}}
